% ---- pg-560 （书页 548） ----

\[
F_{\text{nuc}}=\frac{\text{聚类间核能百分比的均方}}{\text{聚类内核能百分比的均方}}=\frac{3.335}{0.255}=13.1
\]
可见不同聚类的公司之间的核能百分比弥散度很大, 而同一聚类内部的公司其核能百分比的变差很小. 燃料成本和年销售额在区分聚类方面看来也有某种重要性.

回顾五个聚类中的公司可以看到, $K$ 均值法给出的结果一般与平均连接分层法的结果一致(见例 12.9). 具有相同或相容地理位置的公司聚集在一起.
还有, 处于同一聚类中的公司其核能百分比大致相同. \hfill$\blacksquare$

正如我们在全书一直强调的那样, 在聚类过程中评价单个变量的重要性时一定要小心, 这种评价必须从多元的角度判断.
所有变量(多元观测值)决定各个聚类均值和各个项目的归属. 此外, 量度个别变量重要性的描述性统计量的值是聚类个数及最终聚类结构的函数.
另一方面, 在得出结果之后, 有些描述性量度可能有助于评估聚类方法是否“成功”.

\mvsub{对非分层方法的最后评注}

对非分层方法在事先不固定聚类数目 $K$ 的问题上有很多争议. 争议包括以下几点:

\noindent 1.\quad 若两个或更多个种子点碰巧处于同一聚类之中, 所得到的最终聚类会难以区分.

\noindent 2.\quad 当存在某个离群值时, 至少会产生一个所含项目非常分散(不集中)的聚类.

\noindent 3.\quad 即使知道总体由 $K$ 个组组成, 抽样过程也有可能产生这样的结果: 来自最稀疏的那一组的数据未出现在样本之中. 这种情况之下, 强行将数据分成 $K$ 组会造成某些毫无意义的聚类.

\noindent 当某一次算法运行中需要用户指定 $K$ 值时, 使这种算法对若干个不同的初值重复运行多次, 始终是一个不错的想法.

其他关于非分层聚类方法的讨论可从文献[3], [8]和[16]中找到.

\section{基于统计模型的聚类}
\phantomsection\label{sec:12.5}

本章前面讨论了若干常用聚类方法, 包括单连接、完全连接、平均连接、王尔德方法, 以及 $K$ 均值聚类等.
这些方法都是直观上合理的聚类程序, 但是都没有模型来解释观测数据是如何产生的.
通过在聚类方法中引进统计模型, 最大的好处就是可以解释 $(p\times 1)$ 维度量 $\mathbf{x}_j$ 是如何从 $N$ 个对象中产生的.
最常用的模型是假定类 $k$ 中的对象所占的期望比例为 $p_k$, 而相应的度量变量的参数有概率密度函数 $f_k(\mathbf{x})$.
如此, 若有 $K$ 个类, 从单个对象得到的观测向量可认为是从一个混合分布中得到, 其概率密度函数为
\[
f_{\text{Mix}}(\mathbf{x})=\sum_{k=1}^{K}p_k f_k(\mathbf{x})
\]
其中 $p_k\ge 0$ 且 $\sum\limits_{k=1}^{K}p_k=1$. 分布 $f_{\text{Mix}}(\mathbf{x})$ 称为是 $K$ 个分布 $f_k(\mathbf{x})$ 的混合物, $k=1,\cdots,K$.
从这个分布中生成的 $N$ 个观测是从分量分布中得到的观测的混合物.

最常用的混合模型是多元正态分布的混合物, 其中 $f_k(\mathbf{x})$ 是 $N_p(\boldsymbol{\mu}_k,\boldsymbol{\Sigma}_k)$ 的密度函数.
对于观测 $\mathbf{x}$, 正态混合模型的密度函数为

\phantomsection\label{eq:12-17}
\begin{equation}
\begin{aligned}
&f_{\text{Mix}}(\mathbf{x}\mid \boldsymbol{\mu}_1,\boldsymbol{\Sigma}_1,\cdots,\boldsymbol{\mu}_K,\boldsymbol{\Sigma}_K)\\
&\quad=\sum_{k=1}^{K}p_k\frac{1}{(2\pi)^{p/2}|\boldsymbol{\Sigma}_k|^{1/2}}\exp\left(-\frac{1}{2}(\mathbf{x}-\boldsymbol{\mu}_k)'\boldsymbol{\Sigma}_k^{-1}(\mathbf{x}-\boldsymbol{\mu}_k)\right)
\end{aligned} \tag{12-17}
\end{equation}

由此模型产生的类有椭球形状, 在中心附近的集中度最高.
